\Article{26}
Wird die Intensität des Erdmagnetismus $T$ durch eine \textit{Zahl} $k$ ausgedrückt, so liegt dieser Zahl eine bestimmte Einheit $V$ zugrunde, nämlich eine Kraft derselben Art. Ihr Zusammenhang mit den anderen, unmittelbar gegebenen Einheiten ist im Vorangehenden enthalten, allerdings in etwas verwickelter Form. Es lohnt sich daher, diesen Zusammenhang hier erneut darzustellen, damit unmittelbar deutlich wird, wie sich die Zahl $k$ ändert, wenn wir von anderen Grundeinheiten ausgehen.

\SourceLeaf{120}
\expandafter\def\csname HistoricalNoteText13\endcsname{Es wird kaum nötig sein, darauf hinzuweisen, dass die den Buchstaben bisher gegebenen Bedeutungen hier nicht mehr gelten.}
Um die Einheit $V$ festzulegen, mussten wir von der Einheit des freien Magnetismus $M$\HistoricalNote{13} und der Einheit des Abstands $R$ ausgehen. Wir setzten $V$ gleich der Kraft, die $M$ im Abstand $R$ ausübt.

Als Einheit $M$ nahmen wir diejenige Menge magnetischer Flüssigkeit, die auf eine gleich große Menge $M$ im Abstand $R$ eine bewegende Kraft (oder, wenn man will, einen Druck) ausübt, die gleich der als Einheit gewählten Kraft $W$ ist. Sie ist also gleich der Kraft, die die als Einheit gewählte beschleunigende Kraft $A$ auf die als Einheit gewählte Masse $P$ ausübt.

Zur Festlegung der Einheit $A$ gibt es zwei Wege. Man kann sie entweder einer unmittelbar gegebenen gleichartigen Kraft entnehmen, etwa der Schwerkraft am Beobachtungsort, oder aus ihrer Wirkung bei der Bewegung von Körpern bestimmen. Beim zweiten Verfahren, dem wir in unseren Rechnungen gefolgt sind, braucht man zwei neue Einheiten, die Zeiteinheit $S$ und die Geschwindigkeitseinheit $C$. Als Einheit $A$ nimmt man die beschleunigende Kraft, die bei einer Wirkungsdauer $S$ die Geschwindigkeit $C$ hervorbringt. Als diese Geschwindigkeit nimmt man schließlich diejenige der gleichförmigen Bewegung über die Strecke $R$ in der Zeit $S$.

Die Einheit $V$ hängt also offenbar von drei Einheiten ab, entweder von $R,P,A$ oder von $R,P,S$.

Nehmen wir nun an, dass wir die Einheiten $V,R,M,W,A,P,C,S$ durch $V',\allowbreak R',\allowbreak M',\allowbreak W',\allowbreak A',\allowbreak P'.\allowbreak\ C',\allowbreak S'$ ersetzen, die untereinander genauso zusammenhängen wie die bisherigen. In der Einheit $V'$ werde der Erdmagnetismus durch die Zahl $k'$ ausgedrückt. Zu untersuchen ist, wie sie mit $k$ zusammenhängt.

Setzen wir
\[
\begin{aligned}
V&=vV'\\
R&=rR'\\
M&=mM'\\
W&=wW'\\
A&=aA'\\
P&=pP'\\
C&=cC'\\
S&=sS'
\end{aligned}
\]
so sind $v,r,m,w,a,p,c,s$ reine Zahlen, und es gilt
\SourceLeaf{121}
\[
\begin{aligned}
kV&=k'V'\quad\text{oder}\quad kv=k'\\
v&=\frac m{rr}\\
\frac{mm}{rr}&=w=pa\\
a&=\frac cs\\
c&=\frac rs
\end{aligned}
\]
Aus der Verbindung dieser Gleichungen ergibt sich
\[
\begin{aligned}
\text{I.}\qquad k'&=k\BareRoot{\frac p{rss}}\\
\text{II.}\qquad k'&=k\BareRoot{\frac{pa}{rr}}
\end{aligned}
\]
Solange wir bei dem Verfahren bleiben, dem wir in unseren Rechnungen gefolgt sind, müssen wir die erste Formel verwenden. Nehmen wir zum Beispiel Meter und Gramm statt Millimeter und Milligramm als Einheiten, so gilt $r=\frac1{1000}$ und $p=\frac1{1000}$, also $k'=k$. Nehmen wir die Pariser Linie und das Berliner Pfund, so gilt $r=\frac1{2{,}255829}$ und $p=\frac1{467711{,}4}$, also $k'=0{,}002196161k$. Damit ergeben zum Beispiel die Versuche VIII den Wert $T=0{,}0039131$.

Ziehen wir es vor, das andere Verfahren zu wählen und die Schwerkraft als Einheit der beschleunigenden Kräfte zu nehmen, so setzen wir für die Göttinger Sternwarte $a=\frac1{9811{,}63}$. Solange wir Millimeter und Milligramm beibehalten, sind die Zahlen $k$ dann mit $0{,}01009554$ zu multiplizieren. Änderungen dieser Einheiten sind nach Formel II zu behandeln.

